
Ultra Peripheral Processes are processes in which two nuclei pass each other with impact parameter larger than their radii sum. Instead, in these types of collisions the interaction is done via the photon flux surrounding the nuclei due to their acceleration. Using the EPA, it then follows that we can describe UPC as a set of processes in which there is a photon interacting either with the target nucleus or with a photon emitted by the target nucleus. In particular, we are interested in a subset of processes in which the final state is a di-jet.
Under such  constraints, we have a few different processes, which are labeled based on the type of interaction:

\begin{description}
    \item [two-photon interaction ($\gamma\gamma$)] In these interactions photons from both nuclei are interacting with each other. This can happen in 3 distinct topologies: direct-direct, in which there is a pure photon-photon interaction; direct-resolved, in which one photon dissociates into a $q\bar{q}$ pair that interacts with the photon; resolved-resolved, in which both photons dissociate and the remaining quarks interact.
    \item [photo-nuclear ($\gamma$A)] In this interaction the photon interacts with nucleons in the target nucleus. This can either result in nuclear breakup through electromagnetic dissociation, or not.
    \item [photo-diffractive ($\gamma \mathbb{P}$)] In this interaction the photon does not interact with the nucleons directly. Instead, it interacts with the pomeron emitted by the target nucleus. The diffraction can then occur by either nucleus-generated pomeron (coherent diffraction) or a nucleon-generated pomeron (incoherent diffraction).
\end{description}

While these are good definitions for the initial state, things become more complicated as one tries to infer the initial state from the final state kinematics.  Due to final state radiation (FSR) and hadronization, the final states in such collisions can vary greatly. 
